\thispagestyle{empty} % Pas de numéro de page
\setcounter{page}{0}  % Ne pas incrémenter la numérotation

\setcode{utf8}

\begin{center}

\vspace{0.1cm}

\fbox{
\begin{minipage}{0.97\textwidth}
\textbf{Résumé} \\
Ce mémoire présente la conception et le déploiement optimisé d’un core 5G SA basé sur les technologies Cloud-Native. L’objectif principal est d’exploiter les avantages des environnements cloud afin d’assurer un déploiement efficace, hautement disponible et sécurisé du core 5G. En s’appuyant sur des outils open source, nous avons élaboré une solution en plusieurs couches intégrant un pipeline CI/CD, un service mesh pour la gestion fine du trafic, ainsi qu’un système complet d’observabilité. Une attention particulière a été accordée à la sécurisation des communications inter-CNFs, avec une micro-segmentation dynamique garantissant l’application du modèle Zero Trust, ainsi qu’à la centralisation avancée des logs. Les résultats obtenus confirment la viabilité d’un déploiement Cloud-Native robuste, évolutif et conforme aux exigences des réseaux 5G de nouvelle génération.
\\
\textbf{Mots-clés :} 5G, Cloud Native, DevSecOps, CI/CD, Zero Trust, Service Mesh.
\end{minipage}
}

\vspace{0.1cm}

\fbox{
\begin{minipage}{0.97\textwidth}
\textbf{Abstract} \\
This thesis presents the design and optimized deployment of a 5G SA core network based on Cloud-Native technologies. The main objective is to leverage the advantages of cloud environments to ensure an efficient, highly available, and secure deployment of the 5G core. Relying on open-source tools, we developed a multi-layered solution integrating a CI/CD pipeline, a service mesh for fine-grained traffic control, and a comprehensive observability system. Particular attention was paid to securing inter-CNF communications through dynamic micro-segmentation, ensuring enforcement of the Zero Trust model, as well as to advanced log centralization. The results obtained confirm the feasibility of a robust, scalable, and Cloud-Native deployment aligned with next-generation 5G network requirements.
\\
\textbf{Keywords:} 5G, Cloud Native, DevSecOps, CI/CD, Zero Trust, Service Mesh
\end{minipage}
}

\vspace*{0.1cm}

\fbox{
\begin{minipage}{0.97\textwidth}
\begin{arabtext}
\textbf{ملخص} \\
يعرض هذا المشروع تصميم ونشر نواة شبكة الجيل الخامس المستقلة بالاعتماد على تقنيات الحوسبة السحابية الأصلية. يهدف هذا العمل إلى الاستفادة من مزايا البيئات السحابية، مثل القابلية للتوسع، والتوافر العالي، والتشغيل الآلي، لضمان نشر فعّال وآمن وعالي الموثوقية لنواة شبكة الجيل الخامس. تم تطوير حل متعدد الطبقات باستخدام أدوات مفتوحة المصدر، يتضمن خط متكامل للتكامل والنشر المستمر، وشبكة خدمات للتحكم الدقيق في حركة البيانات، بالإضافة إلى نظام مراقبة شامل لجميع المكونات. كما تم التركيز بشكل خاص على تأمين الاتصالات بين وظائف الشبكة السحابية من خلال تطبيق نموذج الثقة المعدومة، بالاعتماد على تقنيات التقسيم الشبكي الديناميكي، إلى جانب نظام مركزي متقدم لتجميع وتحليل السجلات. وتؤكد النتائج المحققة قابلية تنفيذ نشر سحابي قوي، مرن، وقابل للتوسّع، يتماشى مع متطلبات شبكات الجيل الخامس الحديثة.
\\
\textbf{كلمات مفتاحية} :
الجيل الخامس، الحوسبة السحابية، التطوير والتأمين والعمليات، التكامل والنشر المستمر، نموذج الثقة المعدومة.
\end{arabtext}
\end{minipage}
}

\end{center}